\SourceLeaf{230}%
\Article{5}
\indent Die beiden im vorigen Artikel gefundenen Lösungen beziehen sich offenbar auf gegenüberliegende Punkte der Kugeloberfläche oder, anders gesagt, auf entgegengesetzte Richtungen. Das entspricht der Natur der Sache, denn man kann die Normale nach jeder der beiden Seiten der gekrümmten Fläche ziehen. Will man die beiden an die Fläche angrenzenden Bereiche unterscheiden und den einen als Außenbereich, den anderen als Innenbereich bezeichnen, so lässt sich mithilfe des in Artikel 2 (VII) entwickelten Satzes auch jeder Normalen die zugehörige Lösung zuordnen, sobald ein Kriterium zur Unterscheidung der beiden Bereiche feststeht.

Bei der ersten Methode ist ein solches Kriterium aus dem Vorzeichen des Wertes der Größe \(W\) zu gewinnen. Im Allgemeinen trennt die gekrümmte Fläche diejenigen Teile des Raumes, in denen \(W\) einen positiven Wert annimmt, von denen, in denen der Wert von \(W\) negativ wird. Aus jenem Satz folgt leicht, dass man die erste Lösung wählen muss, wenn \(W\) zur Außenseite hin einen positiven Wert annimmt und man sich die Normale nach außen gezogen denkt. In jedem Fall lässt sich leicht entscheiden, ob auf der ganzen Fläche dieselbe Regel für das Vorzeichen von \(W\) gilt oder ob für verschiedene Teile verschiedene Regeln gelten. Solange die Koeffizienten \(P,Q,R\) endliche Werte haben und nicht alle drei zugleich verschwinden, verhindert das Stetigkeitsgesetz jeden Wechsel.

Folgen wir der zweiten Methode, so können wir uns auf der gekrümmten Fläche zwei Kurvenscharen vorstellen. Bei der einen ist \(p\) variabel und \(q\) konstant, bei der anderen ist \(q\) variabel und \(p\) konstant. Die gegenseitige Lage dieser Kurven in Bezug auf die Außenseite muss darüber entscheiden, welche Lösung zu wählen ist. Wir betrachten den Zweig einer Kurve der ersten Schar, der vom Punkt \(A\) bei wachsendem \(p\) ausgeht, den Zweig einer Kurve der zweiten Schar, der vom Punkt \(A\) bei wachsendem \(q\) ausgeht, und die nach außen gezogene Normale. Liegen diese drei Linien \textit{ebenso} wie die jeweiligen Achsen \(x,y,z\) vom Ursprung der Abszissen aus (wenn man sich zum Beispiel in jeder der beiden Dreiergruppen die erste Linie nach links, die zweite nach rechts und die dritte nach oben gerichtet denken kann), so ist die erste Lösung zu wählen. Ist die gegenseitige Lage der drei Linien derjenigen der Achsen \(x,y,z\) entgegengesetzt, so gilt die zweite Lösung.

Bei der dritten Methode ist zu prüfen, ob ein positiver Zuwachs von \(z\) bei unveränderten \(x\) und \(y\) zur Außenseite oder zur Innenseite führt. Für die nach außen gerichtete Normale gilt im ersten Fall die erste Lösung, im zweiten die zweite.
